\bookchapter{Глава 36: Телефон и умный дом — делаем дом умнее}

\chapterintro{У телефона есть ещё одна всё более важная роль — «пульт» всего дома. От лампы и розетки до дверного замка умный дом превращается из «игрушки для гиков» в «домашнюю норму». Эта глава поможет разобраться в протоколах, экосистемах, хабах и стартовых наборах, а в конце — чек-лист безопасности, чтобы «умность» не обернулась проблемами.}

\booksection{Протоколы умного дома: выбор имеет значение}

Перед покупкой умного устройства смотрите, какой протокол оно поддерживает: от этого зависит, сможет ли им управлять ваш телефон и экосистема. Сейчас основных четыре, у каждого свой характер:

\begin{itemize}
  \item Matter: выпущен в 2022 году CSA (Connectivity Standards Alliance), за ним — Apple, Google, Amazon, Samsung и другие гиганты. Цель одна — чтобы устройства всех брендов говорили на одном языке. В 2023–2026 годах быстро становится единым стандартом; поддерживающие устройства могут одновременно работать с Apple HomeKit, Google Home, Amazon Alexa и Mi Home, не привязываясь к одной экосистеме.
  \item Протокол Thread: сетевой «трубопровод» низшего уровня для Matter. Низкое энергопотребление, самоорганизующаяся сеть (если узел обесточен, сигнал автоматически идёт в обход), быстрый отклик без подвисаний. Покупайте устройства с пограничным маршрутизатором Thread (например, HomePod mini, Apple TV 4K, хаб-шлюз Xiaomi) — стабильность будет заметно выше.
  \item Wi-Fi: широкая полоса пропускания и прямое подключение к роутеру, но высокое энергопотребление; при большом числе устройств каналы забиваются. Подходит для камер, роботов-пылесосов и других устройств, передающих большие объёмы данных.
  \item Bluetooth (BLE): простая настройка и экономия энергии, но малая дальность и нет самоорганизующейся сети. Подходит для устройств с малым объёмом данных: замков, термометров и гигрометров.
  \item Zigbee: проверенный низкопотребляющий протокол с самоорганизующейся сетью, требует отдельного шлюза; устройства дёшевы и стабильны, но совместимость между брендами слабая, и его постепенно вытесняет связка Matter+Thread.
\end{itemize}

\begin{tcbinfo}{Главное правило одной фразой}
Покупая умные устройства в 2026 году, в первую очередь смотрите, есть ли маркировка «Matter over Thread» — сейчас это выбор, который лучше всех выдерживает проверку временем. Старые устройства, поддерживающие только Wi-Fi или Bluetooth, не спешите выбрасывать: пользуйтесь ими, а при новых покупках переходите на Matter.

\end{tcbinfo}

\booksection{Сравнение трёх экосистем}

Протокол — это «язык», экосистема — «круг общения». Даже если всё поддерживает Matter, для автоматизаций и голосового управления вам всё равно придётся выбрать одну основную экосистему. Сравнение трёх ведущих экосистем ниже:

\begin{center}
\begin{tabularx}{\linewidth}{|X|X|X|X|}
\hline
\textbf{Экосистема} & \textbf{Сильные стороны} & \textbf{Слабые стороны} & \textbf{Кому подходит} \tabularnewline
\hline
Apple HomeKit & Лучшая приватность (в основном локальная обработка), плавная связка с iPhone/Siri/Shortcuts & Устройств немного и стоят недёшево, нужен хаб — HomePod или Apple TV & Те, у кого вся техника Apple и кто ценит приватность и качество опыта \tabularnewline
Mi Home (Xiaomi) & Самый богатый выбор устройств, лучшая цена-качество, удобный ассистент Xiao AI & Экосистема довольно закрытая, много данных уходит в облако, много рекламы & Начинающие с ограниченным бюджетом, которые хотят обустроить весь дом в одном месте \tabularnewline
Huawei Smart Life (HarmonyOS) & Бесшовная передача задач между телефонами/планшетами/часами HarmonyOS; количество устройств быстро растёт & Немного зарубежных устройств, экосистема ещё догоняет конкурентов & Пользователи Huawei, ценящие межустройственную связку «Супертерминал» \tabularnewline
\hline
\end{tabularx}
\end{center}

\begin{tcbinfo}{Как выбрать экосистему: практические советы}
Выбирайте экосистему по телефону: владельцам iPhone — прежде всего HomeKit, пользователям Huawei — Smart Life, остальным на Android проще всего с Mi Home. Благодаря протоколу Matter устройствами разных экосистем можно управлять с одного телефона, так что «выбирать лагерь» раз и навсегда не обязательно.

\end{tcbinfo}

\booksection{Телефон как центр умного дома}

В большинстве базовых сценариев телефон и есть ваш «центр умного дома» — через приложение вы управляете светом, кондиционером, шторами и безопасностью, а ещё можете настроить автоматизации, чтобы дом «оживал сам»:

\begin{itemize}
  \item Управление устройствами: откройте приложение HomeKit/Mi Home/AI Life — можно удалённо или локально включать и выключать свет, менять цветовую температуру, задавать температуру кондиционера, закрывать и раздвигать шторы, смотреть картинку с камер.
  \item Сценарии автоматизации: комбинируйте «условие действие». Три самых ходовых сценария: Режим «Дома»: геолокация показала, что вы дома + стемнело включить свет в прихожей + включить кондиционер + закрыть шторы + запереть дверь.
  \item Режим «Не дома»: последний вышел из дома (геолокация или дверь заперта) выключить свет + выключить кондиционер + отключить розетки + включить камеры + поставить охрану на сигнализацию.
  \item Ночной режим: каждый день в 23:00 выключить основной свет во всём доме + закрыть шторы + кондиционер на 26 °C + включить ночник.
\end{itemize}

NFC одним касанием: наклейте NFC-метку у двери — одно касание телефоном запускает нужный сценарий (коснулись — включился свет в прихожей, ещё раз — включился ночной режим). На iPhone автоматизация делается через «Быстрые команды» (Shortcuts), на Android метку можно записать через приложение «Кошелёк» или фирменное приложение производителя.

\begin{tcbpractice}{Попробуйте: соберите режим «Дома»}
В приложении умного дома на телефоне создайте автоматизацию: условие — «прибытие в место» (дом), действия по порядку — «включить свет в прихожей», «кондиционер на 26 °C», «закрыть шторы». Сохраните, выйдите погулять и посмотрите, загорится ли свет само, когда вы вернётесь — радость от первого успеха приятнее, чем от заказа нового устройства.

\end{tcbpractice}

\booksection{Рекомендации по стартовому набору}

Не обязательно скупать всё сразу: с этой четвёрки вы закроете около 80\% повседневных нужд:

\begin{itemize}
  \item Умная колонка: Xiao AI (экосистема Mi Home, выгодная цена) / HomePod mini (экосистема Apple, заодно хаб Thread) / Alexa (широкая экосистема за рубежом). Внимание: колонка определяет вашу голосовую экосистему — сначала выберите экосистему, потом покупайте колонку.
  \item Умная лампочка: Yeelight (Mi Home/Matter, регулируются цветовая температура и яркость, недорого) / Philips Hue (отличные цвета, стабильность, но нужен шлюз Hue, дороже). Внимание: перед покупкой проверьте цоколь (E27/B22) и нужна ли вам регулировка цвета.
  \item Умная розетка: умная розетка Xiaomi (дёшево, есть учёт электроэнергии) / TP-Link Kasa (хорошая поддержка Matter за рубежом). Внимание: при использовании с мощной техникой (водонагреватель, обогреватель) обязательно смотрите номинальную мощность, чтобы не перегрузить.
  \item Умный замок: Lockin/Xiaomi (цена-качество, связка с Mi Home) / Kaadas (крупный производитель классических замков, надёжная механика). Внимание: проверьте поддержку Matter или вашей экосистемы, обратите внимание на заряд батареи и расположение резервной замочной скважины.
\end{itemize}

\begin{tcbinfo}{В каком порядке покупать}
Сначала умная колонка (определяет экосистему) затем умные лампочки и розетки (низкий порог входа, результат виден сразу) в конце замки и камеры (это безопасность, выбирайте осторожно). Добавляйте за один раз только один тип устройств, осваивайте и расширяйте дальше — так вы избежите ситуации «купил, но не умею пользоваться».

\end{tcbinfo}

\booksection{Вопросы безопасности}

Чем больше в доме умных устройств, тем больше поверхность атаки. Вот четыре обязательных правила — проверьте себя по каждому:

\begin{enumerate}
  \item Меняйте пароли по умолчанию: заводские пароли камер, роутеров и замков давно опубликованы в интернете — первое, что нужно сделать после покупки, это установить длинный неповторяющийся надёжный пароль.
  \item Вовремя обновляйте прошивки: обновления от производителя часто закрывают уязвимости; включите в приложении «автообновление» или хотя бы раз в месяц проверяйте вручную.
  \item Изоляция сети: разместите IoT-устройства в отдельной Wi-Fi-сети (или гостевой сети), отдельно от основной сети телефона и компьютера. Если одно устройство взломают, добраться до ваших важных данных будет куда труднее.
  \item Защита приватности: включите на камерах маскировку «приватных зон» (спальня и санузел не снимаются), отключите круглосуточную облачную запись; не оставляйте надолго доступ сторонним приложениям к картинке с камер, когда вас нет дома.
\end{enumerate}

\textbf{Вот этих граблей избегайте}\par

\begin{itemize}
  \item Не открывайте незнакомцам дверь через функцию «удалённое открывание замка» — лучше используйте временный пароль, который перестаёт действовать после использования.
  \item Не направляйте камеру прямо на кровать и в окно — вы разглашаете собственную приватность.
  \item Не покупайте дешёвые безымянные устройства — они почти не получают обновлений прошивки, и вы фактически ставите в дом «чёрный ход».
\end{itemize}

\begin{tcbsummary}{Краткие итоги главы}
\begin{itemize}
  \item Matter+Thread — общее направление 2023–2026 годов; покупая новое устройство, в первую очередь смотрите на эту маркировку
  \item Три экосистемы сильны по-своему: Apple делает упор на приватность и опыт, Mi Home — на цену и качество, Huawei — на слаженную работу множества устройств
  \item Телефон может быть центром дома: соберите три сценария «Дома / Не дома / Сон» — и дом начнёт жить сам
  \item Стартовая четвёрка: умная колонка, лампочка, розетка и замок — добавляйте постепенно в этом порядке
  \item Сменить пароль по умолчанию, обновлять прошивку, изолировать сеть, скрывать приватные зоны — четыре обязательных правила безопасности
\end{itemize}

\end{tcbsummary}
